
\documentclass{article}
\usepackage{colortbl}
\usepackage{makecell}
\usepackage{multirow}
\usepackage{supertabular}

\begin{document}

\newcounter{utterance}

\centering \large Interaction Transcript for game `ifeval', experiment `ifeval', episode 54 with playornotplay\_\_qwen3.5{-}2b\_\_playornotplay{-}v1.0{-}merged{-}fp32{-}828e356.
\vspace{24pt}

{ \footnotesize  \setcounter{utterance}{1}
\setlength{\tabcolsep}{0pt}
\begin{supertabular}{c@{$\;$}|p{.15\linewidth}@{}p{.15\linewidth}p{.15\linewidth}p{.15\linewidth}p{.15\linewidth}p{.15\linewidth}}
    \# & \multicolumn{2}{c}{Player} && \multicolumn{2}{c}{Game Master} \\
    \hline

    \theutterance \stepcounter{utterance}  
    & & & \multicolumn{4}{p{0.6\linewidth}}{
        \cellcolor[rgb]{0.9,0.9,0.9}{
            \makecell[{{p{\linewidth}}}]{
                \texttt{\tiny{[P1$\langle$GM]}}
                \texttt{Take the text below as a starting point, and make it a complete article: "You may have to meet with a helper to work out a parenting plan. The first would be to estimate how many time you have everyday for parenting, and is that enough...."} \\
\texttt{Avoid using the following keywords: sleep, cook, feed} \\
\texttt{Mention the keyword 'schedule' for more than 5 times.} \\
            }
        }
    }
     \\ \\

    \theutterance \stepcounter{utterance}  
    & \multicolumn{4}{p{0.6\linewidth}}{
        \cellcolor[rgb]{0.561,0.737,0.561}{
            \makecell[{{p{\linewidth}}}]{
                \texttt{\tiny{[P1$\rangle$GM]}}
                \texttt{You may have to meet with a helper to work out a parenting plan. The first step is to estimate how many time you have everyday for parenting, and is that enough. You will need to create a **schedule** that works for both you and your partner. This **schedule** should include the times you will be available for each other.} \\
\\ 
\texttt{Once you have a rough idea of your daily availability, you can start filling out the **schedule** with specific activities. For example, you might have one hour in the morning, one hour in the afternoon, and one hour in the evening. You can use this **schedule** to decide which activities are most important for each parent.} \\
\\ 
\texttt{The next step is to decide which activities are shared. You can both do some activities, or you can split them. For instance, you might both do the morning and evening activities, while your partner does the afternoon activity. This **schedule** will help you both understand how to divide the responsibilities.} \\
\\ 
\texttt{You can also use the **schedule** to include breaks. You might need to take a break every hour or two to rest or to eat a meal. This **schedule** will help you both know when to take these breaks.} \\
\\ 
\texttt{Finally, you can use the **schedule** to communicate with your partner. You can both review the **schedule** together to make sure it works for both of you. If you need to make changes, you can adjust the **schedule** accordingly.} \\
\\ 
\texttt{By following these steps, you can create a **schedule** that works for both of you. This **schedule** will help you both understand how to divide the responsibilities and ensure that you both get the time you need.} \\
            }
        }
    }
    & & \\ \\

    \theutterance \stepcounter{utterance}  
    & & & \multicolumn{2}{p{0.3\linewidth}}{
        \cellcolor[rgb]{0.9,0.9,0.9}{
            \makecell[{{p{\linewidth}}}]{
                \texttt{\tiny{[GM$|$GM]}}
                {'keywords:forbidden_words': {'forbidden_words': ['sleep', 'cook', 'feed']}, 'keywords:frequency': {'relation': 'at least', 'keyword': 'schedule', 'frequency': 6}}
            }
        }
    }
    & & \\ \\

    \theutterance \stepcounter{utterance}  
    & & & \multicolumn{2}{p{0.3\linewidth}}{
        \cellcolor[rgb]{0.9,0.9,0.9}{
            \makecell[{{p{\linewidth}}}]{
                \texttt{\tiny{[GM$|$GM]}}
                \texttt{game\_result = WIN} \\
            }
        }
    }
    & & \\ \\

\end{supertabular}
}

\end{document}
